\usetikzlibrary{fit}

\tikzset{
  umlactorstyle/.style={
    minimum width=1.2cm,
    minimum height=1.6cm,
    shape=rectangle,
    inner sep=0pt,
    path picture={
      \begin{scope}[shift={(path picture bounding box.center)}, yshift=0.2cm]
        \draw[thick] (0,0.4) circle (0.18);
        \draw[thick] (0,0.22) -- (0,-0.2);
        \draw[thick] (-0.3,0.05) -- (0.3,0.05);
        \draw[thick] (0,-0.2) -- (-0.2,-0.5);
        \draw[thick] (0,-0.2) -- (0.2,-0.5);
      \end{scope}
    }
  }
}

\renewcommand{\umlactor}[2][]{%
  \node [umlactorstyle, label={[align=center]below:#2}, #1] (#2) {};
}

\renewcommand{\umlusecase}[2][]{%
  \node [
    draw, 
    ellipse, 
    align=center, 
    minimum width=3cm, 
    minimum height=1cm, 
    #1
  ] {#2};
}

\RenewDocumentEnvironment{umlsystem}{O{} m +b}{
  \begin{scope}
    #3 
    % Create a coordinate to track the internal bounding box
    \coordinate (temp_boundary) at (current bounding box.center);
    \node (internal_fit) [fit=(current bounding box), inner sep=15pt] {};
  \end{scope}

  \begin{scope}[on background layer]
    % This node is now named #2 (e.g., Сайт)
    \node[
      draw, 
      rectangle, 
      #1,
      fill=black!3,
      fit=(internal_fit),
      label={[anchor=north, font=\bfseries]north:#2}
    ] (#2) {}; % <--- This is the crucial part
  \end{scope}
}{}

\renewcommand{\umlassoc}[2][]{
  \draw [#1] #2;
}

\newcommand{\umlrel}[3][]{
  \draw [->, >=latex, dashed, shorten >=2pt, #1] #2 
    node [pos=0.7, above, font=\footnotesize] {$\ll$#3$\gg$};
}